%% parts/p06/axioms-hierarchy.tex --- Phase-4 P6 foundations/ordinals slice
%%
%% Source accounting: v1.1 lines 10063--10557.  The repaired ordinal
%% trichotomy route (R-034) is preserved in annex:p6.ordinals.1.  The
%% distinct ordered-pair statement and proof at v1.1 10283--10299
%% (source-unit ordinals 710--711) remain owed verbatim in exact source
%% interval SRC-10272-10352.  The related Part-I source presentation at
%% source-unit ordinals 16--17 is different and cannot discharge this one.
%% Depth/apparatus: D-P7-001, N-a-066, A-058, A-060, A-129, A-130.

\chapter{Axioms, Ordinals, and the Cumulative Hierarchy}
\label{sec:p06-axioms-hierarchy}
\chaptermark{Axioms, Ordinals, and Hierarchy}

The entries in this chapter separate set theory as an object theory from the
ambient set theory used to conduct metamathematics.  The source organization
and the cited background are those of \cite{Jech2003,Kunen2011}.

\section[The language and axioms of ZF and ZFC]{The language and axioms of
\(\mathsf{ZF}\) and \(\mathsf{ZFC}\)}
\label{sec:p06-zfc-axioms}

\begin{hbconvention}{def:p06-scoped-semantic-notation}{Scoped semantic notation}
In semantic passages, \(\mathcal M\) and \(\mathcal N\) are the default
metavariables for structures, and \(\mathcal H\) is reserved for an explicitly
declared Henkin structure.  The word \emph{model} means a structure together
with an assertion that it satisfies a specified sentence or theory.  Any
different local use---for example, a forcing name---must be declared at the
start of its scope.
\end{hbconvention}

\begin{hbconvention}{def:p06-membership-language}{Membership language and class notation}
The object language of set theory is the one-sorted first-order language
\(\Lang_{\in}=\{\in\}\), with logical equality and one binary relation symbol.
For a formula \(\varphi(x,\vec p)\), the notation
\(\{x:\varphi(x,\vec p)\}\) denotes a metalinguistically specified
\emph{class}; it introduces neither a set nor a class variable into
\(\Lang_{\in}\).  It denotes a set only when an axiom or theorem supplies a
set with exactly those members.  Thus ``class'' names a metalinguistic
collection here, not a second object sort.
\index{membership language}\index{class}\index{proper class}
\end{hbconvention}

\begin{hbwarning}{def:p06-object-meta-choice}{Two scopes of Choice}
The Axiom of Choice in Table~\ref{tab:p06-zfc-axioms} is a sentence of
\(\Lang_{\in}\) and belongs to \(\mathsf{ZFC}\).  An invocation of Choice in
the surrounding metatheory is instead a principle used to reason about
structures, proofs, or sets from outside the object theory.  Neither use
silently supplies the other.
\end{hbwarning}

The bounded-quantifier notation in Table~\ref{tab:p06-zfc-axioms} is an
object-language abbreviation:
\[
 \forall x\in a\,\varphi\defeq\forall x\,(x\in a\to\varphi),
 \qquad
 \exists x\in a\,\varphi\defeq\exists x\,(x\in a\land\varphi).
\]
Also
\(x\subseteq a\defeq\forall z\,(z\in x\to z\in a)\) and
\(S(x)\defeq x\cup\{x\}\).  The Kuratowski pair
\(\langle a,b\rangle\defeq\{\{a\},\{a,b\}\}\) is the eliminable abbreviation
used below.  The Part~1 presentation of its identity theorem is a separate
inherited source unit; it does not replace or discharge the distinct later
Part~6 statement and proof identified as OPEN below.  For the Choice row, the
following eliminable object-language abbreviations are used:
\[
\begin{aligned}
 \operatorname{Fun}_{\in}(f)\defeq{}&
   \bigl[\forall z\in f\,\exists a\,\exists b\,
          z=\langle a,b\rangle\bigr]\land{}\\
 &\forall a\,\forall b\,\forall c\,
   \bigl((\langle a,b\rangle\in f\land
          \langle a,c\rangle\in f)\to b=c\bigr),\\
 \dom f=A\defeq{}&
   \forall a\,
   \bigl(a\in A\leftrightarrow
         \exists b\,\langle a,b\rangle\in f\bigr).
\end{aligned}
\]
These displays introduce metalinguistic shorthands for formulas, not new
symbols of \(\Lang_{\in}\).  Every schema row below denotes its universal
closure in the displayed parameters \(\vec p\); no parameter remains free in
an axiom instance.  Finally,
\(\exists!y\,\psi(y)\defeq
\exists y\,[\psi(y)\land\forall z\,(\psi(z)\to z=y)]\), with \(z\) fresh.

\begingroup\footnotesize
\begin{longtable}{@{}p{0.16\textwidth}p{0.54\textwidth}p{0.22\textwidth}@{}}
\caption{The axioms and schemas of \(\mathsf{ZFC}\), with their formal side
conditions (apparatus A-058).}\label{tab:p06-zfc-axioms}\\
\toprule
Axiom or schema & Object-language formula & Scope and side conditions \\
\midrule
\endfirsthead
\toprule
Axiom or schema & Object-language formula & Scope and side conditions \\
\midrule
\endhead
Extensionality\index{Axiom of Extensionality} &
\(\forall a\,\forall b\,[\forall x\,(x\in a\leftrightarrow x\in b)\to a=b]\) &
Membership determines set identity. \\
\addlinespace
Empty set\index{empty-set axiom} &
\(\exists a\,\forall x\,(x\notin a)\) &
Extensionality makes the witness unique; it is \(\varnothing\). \\
\addlinespace
Pairing\index{pairing axiom} &
\(\forall a\,\forall b\,\exists c\,\forall x\,[x\in c\leftrightarrow
(x=a\lor x=b)]\) &
The unique witness is \(\{a,b\}\); \(\{a\}=\{a,a\}\). \\
\addlinespace
Union\index{union axiom} &
\(\forall a\,\exists u\,\forall x\,[x\in u\leftrightarrow
\exists y\,(y\in a\land x\in y)]\) &
The unique witness is \(\bigcup a\). \\
\addlinespace
Power set\index{power-set axiom} &
\(\forall a\,\exists p\,\forall x\,[x\in p\leftrightarrow x\subseteq a]\) &
The unique witness is \(\mathcal P(a)\); \(\subseteq\) expands as above. \\
\addlinespace
Infinity\index{axiom of infinity} &
\(\exists I\,[\varnothing\in I\land
\forall x\,(x\in I\to S(x)\in I)]\) &
Asserts an inductive set; leastness of \(\omega\) is Theorem~\ref{thm:p06-omega-least}. \\
\addlinespace
Separation\index{separation schema} &
\(\forall a\,\exists b\,\forall x\,[x\in b\leftrightarrow
(x\in a\land\varphi(x,\vec p))]\) &
One instance for each \(\Lang_{\in}\)-formula \(\varphi\); \(b\) is not free
in \(\varphi\); parameters \(\vec p\) are allowed. \\
\addlinespace
Replacement\index{replacement schema} &
\(\displaystyle
\forall a\bigl[\forall x\in a\,\exists!y\,\varphi(x,y,\vec p)\to
\exists b\,\forall y\,[y\in b\leftrightarrow
\exists x\in a\,\varphi(x,y,\vec p)]\bigr]\) &
One instance for each \(\Lang_{\in}\)-formula \(\varphi\); \(b\) is not free
in \(\varphi\); the antecedent requires functionality on \(a\). \\
\addlinespace
Foundation\index{axiom of foundation} &
\(\forall a\,[a\ne\varnothing\to
\exists x\in a\,\forall y\in x\,(y\notin a)]\) &
Every nonempty set has a member disjoint from it. \\
\addlinespace
Choice\index{Axiom of Choice} &
\(\displaystyle
\forall A\bigl[\forall x\in A\,(x\ne\varnothing)\to
\exists f\,(\operatorname{Fun}_{\in}(f)\land\dom f=A\land
\forall x\in A\,\exists y\,
(\langle x,y\rangle\in f\land y\in x))\bigr]\) &
The displayed \(\operatorname{Fun}_{\in}\), domain, and ordered-pair notation
are eliminable formulas.  Pairwise disjointness is not assumed. \\
\bottomrule
\end{longtable}
\endgroup

\begin{hbdefinition}{def:p06-zf-zfc}{The theories ZF and ZFC}
\(\mathsf{ZF}\) consists of Extensionality, Empty Set, Pairing, Union, Power
Set, Infinity, Foundation, and every instance of Separation and Replacement
in Table~\ref{tab:p06-zfc-axioms}.  The theory \(\mathsf{ZFC}\) is
\(\mathsf{ZF}\) together with the object-language Axiom of Choice.  Some
axiomatizations derive Empty Set from other axioms; the displayed list is the
fixed axiomatization used here.
\index{ZF}\index{ZFC}
\end{hbdefinition}

\begin{hbdefinition}{def:p06-omega}{The least inductive set}
Write
\[
 \operatorname{Ind}(I)\defeq
 \bigl(\varnothing\in I\land
       \forall x\,(x\in I\to S(x)\in I)\bigr).
\]
After Infinity supplies an inductive set \(I_0\), Separation forms
\[
 \omega\defeq
 \{x\in I_0:\forall J\,(\operatorname{Ind}(J)\to x\in J)\}.
\]
Here \(\operatorname{Ind}\) is an eliminable abbreviation for the displayed
\(\Lang_{\in}\)-formula; it does not add a predicate to the language.
\index{inductive set}\index{omega@\(\omega\)}
\end{hbdefinition}

\begin{hbtheorem}{thm:p06-omega-least}{proved=annex:p6.axioms.1}{Existence and leastness of \(\omega\)}
In \(\mathsf{ZF}\), the set \(\omega\) of
Definition~\ref{def:p06-omega} is inductive and is included in every
inductive set.  It is therefore the unique least inductive set.
\end{hbtheorem}

\begin{hbcorollary}{thm:p06-foundation-no-self-member}{proved=annex:p6.axioms.2}{Foundation excludes self-membership}
\(\mathsf{ZF}\) proves \(a\notin a\) for every set \(a\).  The absence of
self-membership alone is not equivalent to Foundation.
\end{hbcorollary}

\begin{hbtheorem}{thm:p06-set-recursion}{sketched=annex:p6.axioms.3}{Set recursion}
Let \(\psi(x,g,y,\vec p)\) be a formula such that for every set \(x\) and
every function \(g\) with domain \(x\), exactly one set \(y\) satisfies
\(\psi(x,g,y,\vec p)\); write that unique value as \(G(x,g)\).  Then there
is a unique class function \(F\) such that
\[
 F(x)=G\bigl(x,F\mathbin{\upharpoonright}x\bigr)
\]
for every set \(x\), with recursion along membership.  For every set \(a\),
\(F\mathbin{\upharpoonright}a\) is a set function.
\end{hbtheorem}

\begin{hbwarning}{def:p06-no-unrestricted-comprehension}{No unrestricted comprehension}
The notation \(\{x:\varphi(x)\}\) does not ensure that the displayed class is
a set.  Separation produces \(\{x\in a:\varphi(x)\}\) only after a bounding
set \(a\) has been supplied.  In particular, \(\{x:x\notin x\}\) cannot be a
set.
\end{hbwarning}

\paragraph{Source disposition.}
\textbf{OPEN:} The distinct Part~6 Kuratowski ordered-pair lemma and proof at
source lines 10283--10299 (source-unit ordinals 710--711) must be restored
verbatim within exact source-mapped interval \texttt{SRC-10272-10352}.  The
related Part~1 presentation at source-unit ordinals 16--17 in
Section~\ref{sec:set-prerequisites-02} is a separate source unit and does not
replace or discharge the Part~6 statement or proof.  The source exercises
asking for a bounded construction of \(a\times b\) and for the exclusion of a
two-cycle likewise remain owed at their exact source-mapped destination.
Warning~\ref{def:p06-no-unrestricted-comprehension}, the Part~1 product entry,
and the Foundation entry above are related cautions only; they do not replace
or discharge those exercises.

\section{Ordinals, induction, and recursion}
\label{sec:p06-ordinals}

\begin{hbconvention}{def:p06-ordinal-standing}{Ordinals, successors, limits, and suprema}
The definition from the Part~1 cardinal--ordinal register remains active: an
\emph{ordinal} is a transitive set strictly well-ordered by membership.  For
ordinals \(\alpha,\beta\),
\[
 \alpha<\beta\defeq\alpha\in\beta,
 \qquad
 \alpha+1\defeq\alpha\cup\{\alpha\}.
\]
An ordinal of the latter form is a \emph{successor ordinal}; a nonzero ordinal
that is not a successor is a \emph{limit ordinal}.  If \(A\) is a set of
ordinals, \(\sup A\defeq\bigcup A\), the least ordinal at least every member
of \(A\).  The notation \(\operatorname{Ord}\) denotes the metalinguistic
class of all ordinals, not a set.
\index{ordinal}\index{successor ordinal}\index{limit ordinal}
\index{supremum of ordinals}
\end{hbconvention}

\begin{hbtheorem}{thm:p06-ordinal-trichotomy}{sketched=annex:p6.ordinals.1}{Trichotomy and members of ordinals}
For ordinals \(\alpha\) and \(\beta\), exactly one of
\[
 \alpha\in\beta,
 \qquad \alpha=\beta,
 \qquad \beta\in\alpha
\]
holds.  Every member of an ordinal is an ordinal.
\end{hbtheorem}

\begin{hbtheorem}{thm:p06-transfinite-induction}{proved=annex:p6.ordinals.2}{Transfinite induction}
Let \(\varphi(\alpha,\vec p)\) be a formula.  If for every ordinal
\(\alpha\),
\[
 \bigl(\forall\beta<\alpha\,\varphi(\beta,\vec p)\bigr)
 \to\varphi(\alpha,\vec p),
\]
then \(\varphi(\alpha,\vec p)\) holds for every ordinal \(\alpha\).
This is a theorem about a progressive property; it does not by itself define a
function.
\end{hbtheorem}

\begin{hbtheorem}{thm:p06-transfinite-recursion}{sketched=annex:p6.ordinals.3}{Transfinite recursion}
Let \(G\) be a definable class operation assigning one set \(G(f)\) to every
function \(f\) whose domain is an ordinal.  There is a unique class function
\(F\) on \(\operatorname{Ord}\) such that
\[
 F(\alpha)=G\bigl(F\mathbin{\upharpoonright}\alpha\bigr)
\]
for every ordinal \(\alpha\), and
\(F\mathbin{\upharpoonright}\gamma\) is a set for every ordinal \(\gamma\).
Unlike transfinite induction, this theorem asserts existence and uniqueness of
values.  Unlike the bounded cited version in Part~1, it states the coherent
class function on all ordinals.
\end{hbtheorem}

Replacement is essential to the general recursion theorem because it collects
the values constructed at all earlier stages of a set-sized initial segment.

\section{Ordinal arithmetic and Cantor normal form}
\label{sec:p06-ordinal-arithmetic}

\begin{hbdefinition}{def:p06-ordinal-arithmetic}{Ordinal addition, multiplication, and exponentiation}
Ordinal arithmetic is defined by transfinite recursion on the \emph{right}
argument.  For a nonzero limit ordinal \(\lambda\),
\[
\begin{array}{lll}
 \alpha+0=\alpha,
 &\alpha+(\beta+1)=(\alpha+\beta)+1,
 &\alpha+\lambda=\sup_{\xi<\lambda}(\alpha+\xi),\\[2pt]
 \alpha\cdot0=0,
 &\alpha\cdot(\beta+1)=\alpha\cdot\beta+\alpha,
 &\alpha\cdot\lambda=\sup_{\xi<\lambda}(\alpha\cdot\xi).
\end{array}
\]
For \(\alpha>0\), exponentiation is
\[
 \alpha^0=1,
 \qquad
 \alpha^{\beta+1}=\alpha^\beta\cdot\alpha,
 \qquad
 \alpha^\lambda=\sup_{\xi<\lambda}\alpha^\xi;
\]
and \(0^0=1\), while \(0^\beta=0\) for \(\beta>0\).  These are ordinal
operations; they are not cardinal addition, multiplication, or exponentiation.
\index{ordinal addition}\index{ordinal multiplication}
\index{ordinal exponentiation}
\end{hbdefinition}

\begin{hbtheorem}{thm:p06-ordinal-arithmetic-laws}{cited=Jech2003}{Basic laws of ordinal arithmetic}
For ordinals \(\alpha,\beta,\gamma\), addition and multiplication are
associative,
\[
 \alpha\cdot(\beta+\gamma)=
 \alpha\cdot\beta+\alpha\cdot\gamma,
 \qquad
 \alpha^{\beta+\gamma}=\alpha^\beta\cdot\alpha^\gamma.
\]
For fixed left argument, all three operations satisfy their defining supremum
clause at nonzero limits (for exponentiation, with nonzero base).  Addition and
multiplication are not commutative, and multiplication need not distribute in
the other direction.  The exact law-and-failure register is
Table~\ref{tab:p06-ordinal-arithmetic}.
\end{hbtheorem}

\begingroup\small
\begin{longtable}{@{}p{0.19\textwidth}p{0.46\textwidth}p{0.27\textwidth}@{}}
\caption{Ordinal-arithmetic laws, continuity clauses, and diagnostic
failures (apparatus A-060).}\label{tab:p06-ordinal-arithmetic}\\
\toprule
Operation & Valid laws and limit behavior & Failures and cautions \\
\midrule
\endfirsthead
\toprule
Operation & Valid laws and limit behavior & Failures and cautions \\
\midrule
\endhead
Addition &
\(0+\alpha=\alpha=\alpha+0\); associativity;
\(\beta<\gamma\) implies \(\alpha+\beta<\alpha+\gamma\);
\(\alpha+\lambda=\sup_{\xi<\lambda}(\alpha+\xi)\). &
Not commutative: \(1+\omega=\omega<\omega+1\). \\
\addlinespace
Multiplication &
\(1\cdot\alpha=\alpha=\alpha\cdot1\); associativity;
\(\alpha(\beta+\gamma)=\alpha\beta+\alpha\gamma\).
If \(\alpha>0\), then
\(\beta<\gamma\) implies \(\alpha\beta<\alpha\gamma\);
\(\alpha\lambda=\sup_{\xi<\lambda}\alpha\xi\). &
Not commutative: \(2\cdot\omega=\omega<\omega\cdot2\).
The reverse distributive law fails:
\((1+1)\omega=\omega\ne\omega+\omega\). \\
\addlinespace
Exponentiation &
For \(\alpha>0\),
\(\alpha^{\beta+\gamma}=\alpha^\beta\alpha^\gamma\) and
\((\alpha^\beta)^\gamma=\alpha^{\beta\gamma}\).
For \(\alpha\geq2\), exponentiation is strictly increasing in the exponent;
for \(\alpha>0\),
\(\alpha^\lambda=\sup_{\xi<\lambda}\alpha^\xi\). &
The base-zero clauses are exceptional.  Cardinal-exponent laws cannot be
substituted for ordinal-exponent laws. \\
\addlinespace
Cantor normal form &
Finite decreasing powers of \(\omega\), with positive finite coefficients on
the right.  Leading exponents expose the comparison and limit structure. &
The order of summands is essential; rearrangement generally changes the
ordinal. \\
\bottomrule
\end{longtable}
\endgroup

\begin{hbtheorem}{thm:p06-cantor-normal-form}{cited=Jech2003}{Cantor normal form}
Every nonzero ordinal \(\alpha\) has a unique expression
\[
 \alpha=
 \omega^{\beta_0}\cdot n_0+
 \omega^{\beta_1}\cdot n_1+\cdots+
 \omega^{\beta_{k-1}}\cdot n_{k-1},
\]
where \(k\geq1\) is finite,
\(\beta_0>\beta_1>\cdots>\beta_{k-1}\), and each
\(n_i\in\N\setminus\{0\}\).  The coefficient is written on the right because
ordinal multiplication is not commutative.  The empty sum represents zero.
This is the \emph{Cantor normal form} of \(\alpha\)
\cite{Cantor1897,Sierpinski1965}.
\index{Cantor normal form}
\end{hbtheorem}

\begin{hbexample}{def:p06-cnf-arithmetic}{Order sensitivity in Cantor normal form}
The leading \(\omega^2\)-term of the right summand absorbs lower-degree terms
on its left:
\[
 (\omega^3\cdot2+\omega\cdot5+7)
 +(\omega^2\cdot4+3)
 =\omega^3\cdot2+\omega^2\cdot4+3.
\]
Also \((\omega+1)\cdot\omega=\omega^2\), whereas
\(\omega\cdot(\omega+1)=\omega^2+\omega\).  These are computations with the
recursive clauses, not cardinal-arithmetic identities.
\end{hbexample}

\begin{hbdefinition}{def:p06-epsilon-zero}{The ordinal epsilon-zero}
An \emph{epsilon ordinal} is an ordinal \(\varepsilon\) satisfying
\(\omega^\varepsilon=\varepsilon\).  The least such ordinal is
\(\varepsilon_0\).  If \(\alpha_0=1\) and
\(\alpha_{n+1}=\omega^{\alpha_n}\), then
\[
 \varepsilon_0=\sup_{n<\omega}\alpha_n.
\]
Every ordinal below \(\varepsilon_0\) has a finite term obtained from Cantor
normal forms with recursively nested exponents, but there is no single finite
tower height sufficient for all of them.  The term names an ordinal fixed
point, not a proof-theoretic notation system; later proof-theory entries state
separately how notation systems represent ordinals below it.
\index{epsilon ordinal}\index{epsilon-zero@\(\varepsilon_0\)}
\end{hbdefinition}

\begin{figure}[htbp]
\centering
\begin{tikzpicture}[
  x=1cm,y=1cm,
  milestone/.style={draw,rounded corners,fill=black!4,inner sep=2.5pt,
                    align=center,font=\scriptsize},
  note/.style={align=center,font=\scriptsize}
]
  \draw[-{Latex[length=2.5mm]},thick] (0,0) -- (11.2,0);
  \foreach \x in {0.4,1.65,2.8,4.0,5.2,6.5,8.05,9.45,10.8}
    \draw (\x,0.10) -- (\x,-0.10);
  \node[milestone,above=3pt] at (0.4,0) {finite\\ordinals};
  \node[milestone,above=3pt] at (1.65,0) {$\omega$};
  \node[milestone,above=3pt] at (2.8,0) {$\omega+1$};
  \node[milestone,above=3pt] at (4.0,0) {$\omega\cdot2$};
  \node[milestone,above=3pt] at (5.2,0) {$\omega^2$};
  \node[milestone,above=3pt] at (6.5,0) {$\omega^\omega$};
  \node[milestone,above=3pt] at (8.05,0)
    {$\omega^{\omega^\omega}$};
  \node[milestone,above=3pt] at (9.45,0) {$\cdots$};
  \node[milestone,above=3pt] at (10.8,0) {$\varepsilon_0$};
  \node[note,below=5pt] at (4.2,0)
    {finite Cantor sums with exponents below the next milestone};
  \node[note,below=18pt] at (9.1,0)
    {finite towers of unbounded height};
  \node[note,below=5pt] at (9.7,0)
    {$\sup(1,\omega,\omega^\omega,\omega^{\omega^\omega},\ldots)$};
\end{tikzpicture}
\caption{Ordinal milestones below \(\varepsilon_0\) (apparatus A-129;
schematic, not to scale).  The ellipsis is cofinal in \(\varepsilon_0\); it
does not assert a uniform finite bound on exponent-tower height.}
\label{fig:p06-ordinals-epsilon-zero}
\end{figure}

\section{The cumulative hierarchy, rank, and collapse}
\label{sec:p06-cumulative-hierarchy}

\begin{hbdefinition}{def:p06-cumulative-hierarchy}{Cumulative hierarchy}
Transfinite recursion defines
\[
 V_0=\varnothing,
 \qquad
 V_{\alpha+1}=\mathcal P(V_\alpha),
 \qquad
 V_\lambda=\bigcup_{\beta<\lambda}V_\beta
 \quad(\lambda\text{ a limit ordinal}).
\]
The notation
\(V=\bigcup_{\alpha\in\operatorname{Ord}}V_\alpha\) is class notation, not a
set-forming union: \(\operatorname{Ord}\) is a proper class.  Thus
\(V_\alpha\) names one set-valued level, while \(V\) abbreviates the class of
all sets.
\index{cumulative hierarchy}\index{V alpha@\(V_\alpha\)}
\end{hbdefinition}

\begin{hbtheorem}{thm:p06-hierarchy}{sketched=annex:p6.hierarchy.1}{Hierarchy theorem}
Assume \(\mathsf{ZF}\).  Then:
\begin{enumerate}
 \item every \(V_\alpha\) is transitive;
 \item if \(\alpha\leq\beta\), then \(V_\alpha\subseteq V_\beta\);
 \item every set belongs to some \(V_\alpha\); and
 \item for every set \(x\), there is a least ordinal \(\alpha\) such that
       \(x\subseteq V_\alpha\).
\end{enumerate}
\end{hbtheorem}

\begin{hbdefinition}{def:p06-rank}{Set-theoretic rank}
The \emph{rank} of a set \(x\) is
\[
 \operatorname{rk}(x)\defeq
 \min\{\alpha:x\subseteq V_\alpha\}.
\]
It is an ordinal attached to a set, not the cardinality of that set.  The first
hierarchy level containing \(x\) as an element is
\(V_{\operatorname{rk}(x)+1}\), not \(V_{\operatorname{rk}(x)}\).
\index{rank!set-theoretic}
\end{hbdefinition}

\begin{hbtheorem}{thm:p06-rank-recursion}{cited=Jech2003}{Rank recursion and membership decrease}
For every set \(x\),
\[
 \operatorname{rk}(x)=
 \sup\{\operatorname{rk}(y)+1:y\in x\}.
\]
Consequently \(y\in x\) implies
\(\operatorname{rk}(y)<\operatorname{rk}(x)\), and
\(\operatorname{rk}(\alpha)=\alpha\) for every ordinal \(\alpha\).
\end{hbtheorem}

\begin{hbexample}{def:p06-first-hierarchy-levels}{The first hierarchy levels}
\[
 V_0=\varnothing,
 \qquad
 V_1=\{\varnothing\},
 \qquad
 V_2=\{\varnothing,\{\varnothing\}\},
\]
and \(V_3=\mathcal P(V_2)\) has four members.  Hence
\(\operatorname{rk}(\varnothing)=0\),
\(\operatorname{rk}(\{\varnothing\})=1\), and
\(\operatorname{rk}(V_2)=2\).
\end{hbexample}

\begin{figure}[htbp]
\centering
\begin{tikzpicture}[
  level/.style={draw,rounded corners,fill=black!4,minimum width=4.4cm,
                minimum height=0.55cm,align=center,font=\small},
  note/.style={align=left,font=\scriptsize}
]
  \draw[-{Latex[length=2.5mm]},thick] (0,0) -- (0,7.35)
    node[above,font=\small]{ordinal index};
  \node[level] (v0) at (3.0,0.35) {$V_0=\varnothing$};
  \node[level] (v1) at (3.0,1.15) {$V_1=\mathcal P(V_0)$};
  \node[level] (v2) at (3.0,1.95) {$V_2=\mathcal P(V_1)$};
  \node[note] at (3.0,2.65) {$\vdots$};
  \node[level] (vw) at (3.0,3.35)
    {$V_\omega=\bigcup_{n<\omega}V_n$};
  \node[level] (vws) at (3.0,4.15) {$V_{\omega+1}=\mathcal P(V_\omega)$};
  \node[note] at (3.0,4.85) {$\vdots$};
  \node[level] (vk) at (3.0,5.55)
    {$V_\kappa=\bigcup_{\alpha<\kappa}V_\alpha$};
  \node[level] (vks) at (3.0,6.35) {$V_{\kappa+1}=\mathcal P(V_\kappa)$};
  \foreach \y/\lab in {0.35/$0$,1.15/$1$,1.95/$2$,3.35/$\omega$,
                         4.15/$\omega+1$,5.55/$\kappa$,6.35/$\kappa+1$}
    {\draw (-0.08,\y)--(0.08,\y) node[left=3pt]{\lab};}
  \node[note,right=7pt of vk.east,text width=4.0cm] {
    Conditional large-cardinal marker: if $\kappa$ is strongly inaccessible,
    then $(V_\kappa,\in)\entails\mathsf{ZFC}$.  A measurable $\kappa$ is, in
    particular, inaccessible.\\[4pt]
    Since $\operatorname{rk}(\kappa)=\kappa$, the ordinal $\kappa$ itself
    first occurs as an element of $V_{\kappa+1}$.};
\end{tikzpicture}
\caption{The cumulative hierarchy and its rank axis (apparatus A-130).  The
large-cardinal annotation is conditional: the figure locates such a cardinal
if one exists and asserts no large-cardinal existence theorem.}
\label{fig:p06-cumulative-hierarchy}
\end{figure}

\begin{hbdefinition}{def:p06-well-founded-extensional-relation}{Well-founded set relations}
Let \(R\subseteq A\times A\).  The relation \(R\) is \emph{well-founded on
\(A\)} iff every nonempty subset \(B\subseteq A\) has an \(R\)-minimal
member:
\[
 \forall B\subseteq A\,
 \bigl[B\ne\varnothing\to
 \exists b\in B\,\forall c\in B\,\neg(cRb)\bigr].
\]
It is \emph{extensional on \(A\)} iff distinct members of \(A\) have distinct
sets of \(R\)-predecessors, equivalently
\[
 \forall a\in A\,\forall b\in A\,
 \bigl[\forall c\in A\,(cRa\leftrightarrow cRb)\to a=b\bigr].
\]
These are properties of the set relation
\((A,R)\); well-foundedness here is not itself the set-theoretic Axiom of
Foundation.
\index{well-founded relation}\index{extensional relation}
\end{hbdefinition}

\begin{hbtheorem}{thm:p06-mostowski-collapse}{sketched=annex:p6.hierarchy.2}{Mostowski collapse}
Let \(A\) be a set and \(R\subseteq A\times A\) be well-founded and
extensional.  There are a unique transitive set \(B\) and a unique isomorphism
\(\pi:(A,R)\to(B,\in)\).  For every \(a\in A\),
\[
 \pi(a)=\{\pi(b):bRa\}.
\]
\end{hbtheorem}
